\documentclass[paper=a5,listof=totoc]{scrartcl}
\usepackage[ngerman]{babel}
\usepackage{xcolor}
\usepackage{hsbcommon}
\usepackage[hidelinks=true]{hyperref}
\usepackage{multicol}

\usepackage[scale=3,colorspec=0.96,text=internal~use~only]{draftwatermark}

\usepackage[headsepline]{scrlayer-scrpage}

\usepackage[utf8]{inputenc}
\usepackage{setspace}
\usepackage{biblatex}
\usepackage{eurosym}
\addbibresource{literature.bib}
\usepackage{pgfplots}
\pgfplotsset{width=7cm,compat=1.18}
\tikzset{->/.style={-Latex}}

%% define the dates
\newcommand{\mysession}[1]{%
  \ifcase#1\relax\or
    April 13th\or
    April 20th\or
    May 4th\or
    May 11th\or
    May 18th\or
    June 1st\or
    June 8th\or
    June 15th\or
    June 22nd\or
    June 26th (Friday)\or
    June 29th\or
    July 6th\or
    \else
  \fi
}

\title{Einführung in die Energiemärkte}
\author{Jan Suchanek}
% \date{Summer Semester \the\year\\\today}

\automark{section}
\pagestyle{scrheadings}
\ihead{\leftmark}
   \chead{}
   \ohead{Einführung in die Energiemärkte}

\setlength{\parskip}{1ex}

\makeindex

\begin{document}

\onehalfspacing

\maketitle

\begin{abstract}
    Dies ist das Whitecover-Skript für den (noch in Planung befindlichen) Kurs. Es gibt die Folien 1:1 wieder. Eine Hand \faHandPointRight\ verweist auf Folienwechsel. Mit der Ausarbeitung des Foliensatzes wächst dieses Skript mit. Es dient zugleich als Basis für die im Multiple-Choice format online durchgeführte Abschlussprüfung.
\end{abstract}

\newpage

\tableofcontents

\clearpage

%%% Inhalt

\section{Einführung}

\begin{frame}{Kennzahlen Strommarkt 2025}
    \mytext{
    \begin{itemize}
        \item Erzeugung pro Jahr ca. 438,2 Milliarden kWh = TWh
        \item Durchschnittspreis 89,32 Euro / MWh
        \item Erzeugungswert folglich grob 39 Mrd. Euro / Jahr
        \item ca. 3/4 davon OTC und 1/4 börslich gehandelt
        \item Handelsvolumen sehr viel größer (Mehrfachverkäufe)
        \item Summe der Verbraucherzahlungen ebenfalls größer
        \item Preisschwankungen erheblich
        \item 4 Übertragungsnetzebetreiber
        \item 1 Gebotszone, 1 Börse (EEX/EPEX)
        \item Liberalisierter Markt (Unbundling)
        \item Verbindung zu anderen Energiemärkten
    \end{itemize}
    }
    \mycomment{
    
    Zum Vergleich: 

    \begin{itemize}
        \item der Lebensmittelmarkt in Deutschland wird auf ca. 218 Mrd. Euro Umsatz pro Jahr geschätzt
        \item Deutsche geben ca. 2-3\% des Haushaltseinkommens für Strom aus
        \item[] - einkommensschwache allerdings stets mehr als einkommensstarke
    \end{itemize}
    }
\end{frame}


\section{Rollen im Strommarkt}


\begin{frame}{Produzenten, Verbraucher, Verteilnetzbetreiber}
    \mycomment{
    Energienetze verbinden Produzenten von Energie mit Verbrauchern, und Energiemärkte verbinden Verkäufer mit Kunden, Anbietern mit Nachfragern.

    Produzenten sind große Kraftwerksbetreiber, aber auch vergleichsweise kleine Produzenten von Wind- und Solarstrom.

    Verbraucher sind Haushalte und Unternehmen. 

    Die physische Verbindung geschieht über die Stromnetze. Produzenten wie Verbraucher hängen dabei an regionalen Mittelspannungs-Verteilnetzen. 

    Die Übertragung von Strom ist sehr effizient. Immer müssen Last und Erzeugung übereinstimmen, was lokal kaum möglich ist. Daher braucht es eine weitere Netzebene.
    
    Physikalisch \enquote{fließt} wird der Strom allerdings nicht und wird auch nicht \enquote{verbrauch}, sondern er fließt in beide Richtungen.
    }
    \showmytikz[1]{../shared/tikz/tikz-energiemarkt}
    
\end{frame}

\begin{frame}{Übertragungsnetzbetreiber}
    \mycomment{

    Die Verteilnetze sind untereinander durch das Übertragungsnetz auf Hochspannungsebene verbunden.

    Übertragungsnetzbetreiber (VNBs) sind in Deutschland je nach Region 50Hertz, Amprion, TenneT TransnetBW.

    Gesetzlich ist vorgeschrieben, dass ÜNB und VNB getrennte Unternehmen sein müssen. Das ist eine Form von Unbundling.

    Sie betreiben das physische Netz, dürfen aber nicht selbst mit Strom handeln. Das ist eine weitere Form von Unbundling.
    }
    \showmytikz[1,2]{../shared/tikz/tikz-energiemarkt}
    
\end{frame}

\begin{frame}{Lieferanten: Versorger und Händler}
    \mycomment{
    Stromhändler  kaufen Leistung bei Produzenten ein und verkaufen sie an die Lieferanten. Diese versorgen die Endkunden, kümmern sich also um Kundengewinnung, Service und Abrechnung.

    Der gesamte monetäre Layer ist unabhängig vom Netzbetrieb. Die Energieversorger haben mit der technischen Zurverfügungstellung nichts zu tun. 

    Auf den Energiemärkten muss das Angebot immer der Nachfrage entsprechen: es darf nicht mehr Strom verkauft werden, als produziert wird, und nicht mehr eingespeist werden, als abgenommen wird. Das muss sehr genau -- auf die Viertelstunde genau -- funktionieren, denn es gibt kaum Speicher.

    }
    \showmytikz[1,2,3]{../shared/tikz/tikz-energiemarkt}
    
\end{frame}

\begin{frame}{Bilanzkreismanagement}
    \mycomment{
    Daher muss jeder Händler bzw. Versorger genau so viel Strom einzukaufen, wie er verkauft. Dabei gibt es die Möglichkeit, Verantwortung für einen Bilanzkreis zu übernehmen.
    
    Jeder Einspeise- und Entnahmepunkt muss zu einem solchen Bilanzkreis gehören. Dabei geht es um die Marktlokation, oft gleich Messlokation. Es sind also marktliche, nicht physikalische Knoten.
    
    Die Rolle Bilanzkreisverantwortlicher (BKV) bietet sich besonders an, wenn der Händler Kraftwerke oder große Lasten auch steuern bzw. regeln kann. 
    
    Zu dem Bilanzkreismanagement gehört ein Fahrplan, um diesen Ausgleich in jeder Viertelstunde sicherzustellen. 
    
    Die Teilnehmer dieses Bilanzkreises können phyisch weit voneinander entfernt sein. 
    %
    }
    \showmytikz[1,2,3,4]{../shared/tikz/tikz-energiemarkt}
    
\end{frame}

\begin{frame}{Die Strombörse}
    \mycomment{
    Neben direkten Verträgen (Over the Counter, OTC) können Händler und große Produzenten auch die Strombörse benutzen - die EEX oder EPEX spot. Bilanzkreisverantwortliche müssen dies sogar, um ihren Bilanzkreis auszugleichen.
    
    Dort sorgt schon das Marktdesign dafür, dass Angebot und Nachfrage stets übereinstimmen. 
    
    Einige wenige sehr große Verbraucher kaufen ihren Strom direkt an der Strombörse.
    
    Neben direktem Energiekauf (Spotmarkt) werden an der Börse auch Termingeschäfte (Futures) und Optionen (Put- und Calloptionen) gehandelt. Über die Börse sind die Märkte auch Europa- und sogar weltweit verbunden.
    }
    \showmytikz[1,2,3,5]{../shared/tikz/tikz-energiemarkt}
    
\end{frame}

\begin{frame}{Ausgleichsenergie}
    \mycomment{
    Wenn ein Händler bzw. Versorger nicht genug oder zu viel Strom für seinen Bilanzkreis eingekauft hat, so muss er die fehlende Energie  (Ausgleichsenergie) von seinem ÜNB beziehen bzw. abgeben. Der Preis dafür wird von den ÜNB ermittelt. Es wird sichergestellt, dass ein ordentlicher Ausgleich stets lukrativer ist.

    Die Rollen von (Groß-)Händlern und Lieferanten (Stromversorgern) sind nicht notwendigerweise getrennt. Größere Versorger wie Stadtwerke sind häufig auch Bilanzkreisverantwortliche. Gerade kleinere und unabhängige Versorger agieren arbeitsteilig mit Großhändlern, die dann Bilanzkreisverantwortliche sind.

    Die Differenz zwischen günstigem Einkauf an der Strombörse und dem höheren Preis der Ausgleichsenergie ist der Anreiz, Bilanzkreisverantwortung zu übernehmen.
    
    }
    \showmytikz[2,3,4,5,6]{../shared/tikz/tikz-energiemarkt}
    
\end{frame}

\begin{frame}{Netzentgelte}
    \mycomment{
    Verbraucher zahlen darüber hinaus Netzentgelte direkt an die Verteilnetzbetreiber. Diese zahlen ihrerseits Entgelte an die ÜNB.

    Für den Anschluss von Produzenten erhalten die VNB keine Netzentgelte -- allerdings gibt es Ausnahmen bzw. Ausgleichszahlungen bei besonders unausgeglichen Verhältnissen, wenn beispielsweise besonders viele Erneuerbare an einem VNB hängen, der wenige Verbraucher hat.

    Netzentgelte werden erhoben nach einer \enquote{Postage Stamp} Methode und sind unabhängig von der Entfernung, grade wie Briefe bei der Post. Auf der Erzeugerseite sind die Netzbetreiber verpflichtet, Produzenten einen Anschluss anzubieten (Netzausbau). 
    }
    \showmytikz[2,3,6,7]{../shared/tikz/tikz-energiemarkt}
    
\end{frame}

\begin{frame}{Steuern und Umlagen}
    \mycomment{
    Der Staat kontrolliert das gesamte System über Regulierungen, über Steuern und Umlagen und Konzessionen. Die Politik und Höhe der Zahlungen ändern sich vergleichsweise häufig, machen aber einen nicht unbeträchtlichen Teil der Stromrechnung der Verbraucher aus und bestimmen das Marktgeschehen damit deutlich.

    Nicht dargestellt ist hier die Marktrolle der Messstellenbetreiber. Wettbewerbliche Messestellenbetreiber sind unabhängige Unternehmen, die Stromzähler installieren und betreiben. Häufig ist der lokale Verteilnetzbetreiber zugleich Messstellenbetreiber als Grundzuständiger. 
    
    Wir sehen: die Strommärkte sind sehr komplex aufgebaut, und es gibt sehr unterschiedliche Rollen, die häufig von verschiedenen Personen wahrgenommen werden. Oft ist diese Trennung eine juristische Notwendigkeit (Unbundling).
    }
    \showmytikz[1,2,3,4,5,6,7,8]{../shared/tikz/tikz-energiemarkt}
    
\end{frame}

\begin{frame}{Weitere Akteure und Themen}
    \mytext{
    \begin{itemize}
        \item Rollenüberschneidung: \enquote{Prosumer}
        \item kritische Infrastruktur: jeder hängt am Netz
        \item sekundär: Erbauer von Kraftwerken und Netzen
        \item Arbeitsplätze: direkt und indirekt
        \item Preise: gesamtwirtschaftliche Bedeutung
        \item Europäisches Verbundnetz / Europarecht
        \item fossile Brennstoffe: Mobilität, e-Mobilität
        \item Effizienz: Dämmung, Heizung: Gebäudesektor
        \item internationale Konkurrenz: Industriestrompreise
        \item Data Center, KI, Digitalisierung
    \end{itemize}
    }
\end{frame}

\section{Produzenten}

\begin{frame}{Produzenten: Thermische Kraftwerke}
    \mycomment{
    Bei thermischen Kraftwerke treiben Brennstoffe Turbinen und diese wiederum Generatoren an. Sie brauchen Kühlung, und sie erzeugen immer auch Wärme, die entweder als Abwärme verloren geht oder (zum Teil) genutzt wird, beispielsweise als Fernwärme (Kraft-Wärme-Kopplung).

    Diese Turbinen erzeugen direkt die benötigte Frequenz 50 Hertz.

    Es ist möglich, mit der entstehenden Wärme sofort einen weiteren Generator zu betreiben, der ebenfalls Strom erzeugt (GuD-Kraftwerk, Gas- und Dampfturbinenprozess).

    Der produzierte Strom kann natürlich auch andernorts zur Heizung verwendet werden, aber auch zur Kühlung (Klimaanlagen). 
    }
    \showmytikz[1]{../shared/tikz/tikz-energie-kws}
\end{frame}

\begin{frame}{Kinetische Energie}
    \mycomment{
    Fossile (und organische) Brennstoffe können direkt genutzt werden, um Verbrennermotoren zu betreiben. Strom kann genutzt werden, um Elektromotoren zu betreiben. (Eine weitere Nutzung ist natürlich Licht und Elektronik, die allerdings kaum \enquote{zur Last} fällt.)
    }
    \showmytikz[1,2]{../shared/tikz/tikz-energie-kws}
\end{frame}

\begin{frame}{Produzenten: Erneuerbare Energiequellen}
    \mycomment{
    Erneuerbare erzeugen Strom durch Direktumwandlung. Wind und Wasser erzeugen mit Turbinen Wechselstrom; bei Wind wird dann ein Gleichrichter benötigt. Wasserkraftwerke können mitunter direkt 50 Hertz liefern.     
    Solarzellen erzeugen Gleichstrom, der umgewandelt werden muss. 
    
    }
    \showmytikz[1,2,3]{../shared/tikz/tikz-energie-kws}
\end{frame}

\begin{frame}{Neue Entwicklungen: Sektorkopplung}
    \mycomment{
    Strom kann auch genutzt werden, um Wärme zu erzeugen. Auch kann Sonnenergie in Form von Wärme gespeichert werden. Jede solche Umwandlung ist allerdings mit Verlusten verbunden.
    
    }
    \showmytikz[1,2,3,4]{../shared/tikz/tikz-energie-kws}
\end{frame}

\begin{frame}{Neue Entwicklungen: Prosumer}
    \mycomment{
    Die Energiemärkte sind im Wandel, und das nicht nur wegen steigender Anteile der Erneuerbaren. Die Lasten verändern sich und werden zunehmend \enquote{intelligent}. Die Netze müssen an das veränderte Nutzerverhalten angepasst werden - ebenso wie Verträge.
    }
    \mytext{
    \begin{itemize}
        \item die klassische Unterscheidung zwischen Leistung und Last (Produzenten und Verbrauchern) verschwimmt
        \item Einspeisung durch z.B. Balkonkraftwerke, Rückeinspeisung von Electric Vehicles
        \item Energiegemeinschaften, z.B. ganze Dörfer, mit Lasten und Erzeugung
        \item Smart Meter: ggf. Möglichkeiten zur Abregelung
        \item Variable vs. fixeStromtarife
    \end{itemize}
    }
\end{frame}

\section{Preisfindung auf dem Strommarkt}

\begin{frame}{Grundsätzliches}
    \mytext{
    \begin{itemize}
        \item Liberalisierte Strommärkte: 
        \item[] jeder darf verkaufen und kaufen
        \item OTC Over The Counter
        \item börslicher Handel: EEX/EPEX
        \item Spotmärkte, Terminmärkte, Derivate
        \item Händler als eigene Marktrolle (Dienstleister)
        \item Bilanzkreisverantwortliche Händler
        \item Zertifikatehandel (\ce{CO2} etc.)
        \item Verbindung zu internationalen Energiemärkten (fossil)
    \end{itemize}
    }
\end{frame}

% \begin{frame}{Ökonomische Grundlagen}

%     {\color{red}
%     (Umfang und Vorwissen klären, wird noch eingestellt)
%     }
% \end{frame}

\begin{frame}{Merit Order Effekt}
    \mycomment{
    Kraftwerksbetreiber akzeptieren jeden Preis, der oberhalb ihrer Grenzkosten liegt. Diese sind bei Erneuerbaren besonders niedrig. Ordnet man diese Gebote von links nach rechts, so bekommt man die \enquote{Merit Order} genannte Reihenfolge.

    In dieser Darstellung übersteigt die Nachfrage das Angebot von Strom durch Erneuerbare. Verbraucher sind bereit, mehr zu zahlen, nämlich den Preis des Stromangebots durch Kohlekraftwerke. 
    }
    \showmytikz[1]{../shared/tikz/tikz-meritorder-de}
\end{frame}

\begin{frame}{Merit Order Effekt}
    \mycomment{
    Da dieser Preis höher liegt als die Grenzkosten der Erneuerbaren, machen die Erneuerbaren gute Gewinne. Dies ist eine Form der Knappheitsrente.
    }
    \showmytikz[1,2]{../shared/tikz/tikz-meritorder-de}
\end{frame}

\begin{frame}{Merit Order Effekt}
    \mycomment{
    Steigt das Angebot an günstiger Energie, so fällt irgendwann der Preis. 
    
    Da die Einspeisung aus Erneuerbaren stark von der Saison, Tageszeit und Wetterbedingungen abhängt, schwankt die von ihnen angebotene Menge.
    
    Da die Angebotskurve nicht stetig ist, sondern die Merit Order einer Treppe gleicht, kommt es zu starken Preisschwankungen bzw. -sprüngen.

    }
    \showmytikz[3]{../shared/tikz/tikz-meritorder-de}
\end{frame}




\clearpage

\listoffigures
 
\listoftables

\newpage

   \nocite{*}
    \printbibliography[heading=bibintoc]

    \clearpage
    \thispagestyle{empty}
    
    ~
    
    % \cleardoublepage
    % \phantomsection
    % \addcontentsline{toc}{section}{Control Questions}
    % \includepdf[pages=-]{Macroeconomics/questionbank/ME_QuestionPool.pdf}

\cleardoublepage
\phantomsection
\addcontentsline{toc}{section}{Index}
\printindex

\end{document}
